Das Seil ist die Hypotenuse eines rechtwinkligen Dreiecks: $\ell = \lS$ lang, $h = \hS$ hoch. Für den Winkel $\alpha$ des Seils gegen die Waagrechte gilt
\begin{align}
 \sin\alpha &= \saF = \saP{0}{2}, \qquad \cos\alpha = \caF = \caP{0}{2}
\end{align}
Die Zugkraft $F$ hat einen waagrechten Teil $F\cos\alpha$ nach vorn und einen senkrechten Teil $F\sin\alpha$ nach oben.

\begin{abcliste}
\abc
Senkrecht tragen der Boden und der nach oben gerichtete Teil der Zugkraft die Gewichtskraft; waagrecht (konstante Geschwindigkeit) hält der Teil nach vorn der Reibung das Gleichgewicht:
\begin{align}
 \ssc{F}{N} + F\sin\alpha &= m\,g \\
 F\cos\alpha &= \mu\,\ssc{F}{N}
\end{align}
Setzt man $\ssc{F}{N} = m\,g - F\sin\alpha$ in die zweite Gleichung ein, folgt $F\cos\alpha = \mu\,(m\,g - F\sin\alpha)$ und damit
\begin{align}
 F &= \FZF \\
   &= \frac{\muG\cdot\m\cdot\g}{\caP{0}{2}+\muG\cdot\saP{0}{2}} \\
   &= \FZ \approx \result{\FZP{0}{2}}
\end{align}

\abc
\begin{align}
 \ssc{F}{N} &= \FNF \\
   &= \m\cdot\g-\FZ\cdot\saP{0}{2} \\
   &= \FN \approx \result{\FNP{0}{2}}
\end{align}
Leicht nach oben zu ziehen hilft: Es drückt den Schlitten weniger auf den Schnee, und die Reibung ist kleiner.
\end{abcliste}
